\chapter{结论与展望}

\section{结论}
本文围绕“基于 3DGS 与 StyleGAN2 的 MRI 医学图像拟生成平台的设计与实现”开展研究，完成了从 MRI 数据预处理、二维切片生成、连续切片组织、点云初始化、3DGS 训练接入到前后端平台展示的整体设计与实现。论文工作不是单一算法复现，而是面向毕业设计场景，将医学图像数据处理、生成模型训练、三维表达和 Web 平台交互组织为一条可运行的工程链路。该链路能够支撑 MRI 切片生成、实验结果查看、三维资产展示和教学演示等任务。

本文的主要成果体现在四个方面。其一，系统实现了面向 NIfTI 格式 MRI 数据的预处理流程，包括切片提取、灰度归一化、尺寸调整、低质量切片过滤和训练数据整理，为后续模型训练提供了统一输入。其二，系统基于 StyleGAN2-ADA 完成二维 MRI 切片生成模型训练，并将训练权重接入平台推理流程，使用户能够通过随机种子和截断参数生成可复现的 MRI 切片结果。其三，系统将连续切片序列转换为 PLY 初始点云，并构建规则相机参数和训练输入，使 3DGS 能够用于 MRI 切片序列的三维表达实验。其四，系统采用 FastAPI 和 React 实现前后端分离平台，整合结果概览、切片生成、医学影像查看、任务管理、三维展示和 AI 助手等功能，提升了算法实验过程的可视化和可操作性。

本文的创新点主要体现在工程集成和医学图像生成流程适配上。第一，论文将 StyleGAN2 的二维切片生成能力与 3DGS 的三维高斯表达思路结合起来，形成“二维生成结果到三维表达”的平台化流程。第二，论文针对 MRI 切片序列具有规则层间关系这一特点，设计了由连续切片生成初始点云并组织 3DGS 输入的实现方法。第三，论文将训练、推理、结果管理和可视化统一到 FastAPI + React 平台中，使算法结果不再只停留在命令行和离线文件，而能够通过交互界面进行展示和分析。

实验与测试结果表明，平台具备较完整的教学演示和实验复现能力。StyleGAN2 训练链路已经形成可加载权重，平台能够基于该权重生成 MRI 切片。3DGS 部分能够在连续切片监督下完成训练，并输出三维高斯模型、渲染结果和训练指标。第 30000 轮训练达到 L1 为 0.0709、PSNR 为 17.5894 dB，180 对渲染切片与参考切片的平均全局 SSIM 为 0.8935。上述结果说明，本文构建的平台能够把二维生成、三维表达和结果展示连接起来，具有一定的工程应用价值和教学演示价值。

需要说明的是，StyleGAN2、StyleGAN2-ADA 和 3D Gaussian Splatting 的基础理论与核心模型来自已有公开研究，本文不将这些基础算法本身作为原创贡献。本文的本人研究工作主要集中在 MRI 数据处理流程设计、生成模型训练与推理接入、连续切片到点云和 3DGS 输入的工程适配、前后端平台实现，以及实验结果的整理与分析。这样的边界划分有助于明确本文贡献，也避免将已有科研成果误认为本文原创成果。

本研究仍存在一定局限。当前平台主要面向教学演示和工程验证，生成结果不能直接用于临床诊断。实验数据以脑部 T1 MRI 为主，数据类型和病例来源仍有限。3DGS 部分虽然能够完成训练和展示，但 MRI 切片序列不同于普通多视角照片，层间连续性、相机模型、空间尺度和医学结构一致性仍有改进空间。评价指标方面，本文已经使用 L1、PSNR、SSIM 和可视化对比进行分析，但还缺少医学专家评价和下游任务验证。

\section{展望}
后续研究可以从三维表达、评价体系、平台部署和应用场景四个方向继续推进。在三维表达方面，可以进一步结合 MRI 切片的真实层厚、体素间距和空间方向信息，改进相机参数构建方式，并加入层间连续性约束、结构平滑约束和高斯点裁剪策略，以提升边界细节和三维结构稳定性。对于 3DGS 与医学切片序列的结合，还可以探索更适合体数据的采样方式和渲染监督方式，使三维表达更加贴近医学图像本身的空间属性。

评价体系仍需要继续完善。后续可以在现有 L1、PSNR 和 SSIM 的基础上，引入 FID、LPIPS、分割一致性、病灶区域保持能力或下游任务性能等指标。对于医学图像生成任务，仅依靠视觉相似度并不足以说明结果可靠，因此还可以加入医学专业人员的主观评价和错误案例分析，从结构合理性、解剖一致性和教学可用性等角度评价生成结果。

平台工程也有进一步完善空间。当前系统已经能够通过本地 FastAPI 服务和 React 前端完成主要功能展示，但运行过程仍依赖本地模型路径、数据路径、GPU 环境和外部接口配置。后续可以通过配置文件、环境变量、任务队列和容器化部署方式规范运行环境，使平台更适合课程教学、实验复现和多人协作。对于耗时较长的训练任务，也可以加入更完善的日志追踪、失败恢复和资源监控功能。

应用场景方面，本文主要围绕脑部 T1 MRI 切片展开。后续可以扩展到多序列 MRI、CT 或其他公开医学图像数据集，也可以研究跨模态生成、数据增强和三维教学展示等任务。若未来能够获得更充分的数据、专家评价和伦理审查支持，平台还可以作为医学图像生成方法比较、三维重建流程教学和医学人工智能实验课程的辅助工具。
